\section{Conclusion and Limitations}
\label{sec:conclusion}

This draft studies the tokenizer stage of MiniOneRec and argues that graph-structured collaboration should enter semantic ID learning as level-specific structural supervision rather than uniform feature fusion. The resulting method, MGR-SID v1, keeps the semantic \rqvae{} backbone intact, builds a train-only graph bank with coarse, middle, and local views, and applies hierarchy-aware graph regularization to different SID levels. Under a fair upstream-aligned tokenizer regime, this design improves the final generated Industrial SID over the semantic baseline and reduces same-$l_2$ local ambiguity.

The current evidence is still limited in three ways. First, the full training-time result is currently based on Industrial only. Second, the current v1 uses a fixed level-to-view assignment rather than a learned allocation module. Third, we have not yet carried the improved tokenizer into SFT, RL, and end-to-end recommendation evaluation. These limits are intentional: the goal of the present draft is to establish a clean tokenizer-side story before enlarging the experimental scope.

The next step is therefore straightforward. We will use the current best hierarchy-aware tokenizer as the new SID backbone for downstream SFT experiments, and test whether lower local ambiguity at the tokenizer stage translates into fewer same-prefix recommendation errors and better end-to-end recommendation quality.

\paragraph{Reproducibility note.}
All experiments in this draft are built as incremental modules on top of the existing OneRec repository. The default baseline path remains unchanged, and the graph-aware tokenizer branch uses isolated configs, scripts, logs, and outputs.
